\documentclass[letterpaper,11pt]{article}

\usepackage[T1]{fontenc}
\usepackage[utf8]{inputenc}
\usepackage[bitstream-charter]{mathdesign}
\usepackage[top=0.7in,bottom=0.7in,left=0.9in,right=0.9in]{geometry}
\usepackage[hidelinks]{hyperref}
\usepackage{parskip}
\linespread{1.04}
\setlength{\parskip}{7pt}
\pagestyle{empty}
\raggedright

\begin{document}

%----------HEADER----------
\begin{center}
    {\fontsize{18}{20}\selectfont\bfseries DEVA ANAND}\\ \vspace{3pt}
    {\small
    \href{mailto:devaanand@umass.edu}{devaanand@umass.edu}
    \enspace$\vert$\enspace
    (413) 315-1715
    \enspace$\vert$\enspace
    \href{https://devaanand.com}{devaanand.com}
    \enspace$\vert$\enspace
    \href{https://github.com/Deva-1903}{github.com/Deva-1903}}
\end{center}
\vspace{6pt}

\noindent July 15, 2026

\vspace{2pt}

\noindent Old Mission\\
Chicago, IL

\vspace{2pt}

\noindent \textbf{Re: Software Developer}

\vspace{4pt}

\noindent Dear Old Mission Hiring Team,

The part of engineering I enjoy most is making code fast and correct at a level most teams never look at, so a development program built around C++ trading tools, where performance and reliability genuinely matter, is exactly where I want to be. I am completing my MS in Computer Science at UMass Amherst, and I would bring a systems-and-performance habit of mind together with real production backend experience.

My favorite recent project is a from-scratch exercise in getting C++ to the metal: I optimized matrix multiplication from a naive baseline through cache-aware tiling, hand-written AVX/SIMD vectorization with register-aware micro-kernels, and OpenMP multithreading, profiling L1/L3/TLB misses with Linux perf and reading the generated assembly to confirm the compiler did what I expected. The instinct there, measure first and then remove only the bottleneck you can prove, is what I would bring to trading applications. I carried the same habit into data systems, cutting a Spark ETL pipeline's runtime about 25\% on a 1.4-billion-row dataset by tracking shuffle and skew bottlenecks through stage-level timings.

I am not only a systems person, though. Before grad school I spent two years as a backend engineer at Wysa, where I owned a multi-tenant service for 8+ UK clinical clients, diagnosed production incidents across tenants, and shipped fixes under real stakes. Earlier, at Cario, I designed 20+ production REST APIs and migrated a codebase from JavaScript to TypeScript. Across those roles and my projects I have worked with Python, C++, SQL, JavaScript, React, and Linux, on a solid systems-programming and algorithms foundation, which lines up closely with your stack, and I am comfortable jumping into an unfamiliar codebase and becoming useful quickly.

What draws me to Old Mission specifically is the culture you describe: naturally curious people in a team environment who are never quite satisfied with the current version. That is how I work, and a firm that invests in its own engineers rather than chasing outside capital is a place I would give my best. I would love the chance to contribute immediately and to learn the craft of electronic trading from your team. Thank you for considering my application, and I would welcome the chance to talk.

\vspace{6pt}

\noindent Sincerely,\\
Deva Anand

\end{document}
